% ============================================================
% Configuración general
% ============================================================

\usepackage[T1]{fontenc}
\usepackage[utf8]{inputenc}
\usepackage[spanish,es-nodecimaldot]{babel}

% Tipografía académica: serif para cuerpo y sans serif para jerarquía
\usepackage{newtxtext,newtxmath}
\usepackage[scaled=0.92]{helvet}
\usepackage{microtype}

\usepackage[
    top=2.55cm,
    bottom=2.55cm,
    left=3.00cm,
    right=2.55cm,
    headheight=15pt,
    headsep=0.75cm,
    footskip=1.20cm
]{geometry}

\usepackage{setspace}
\setstretch{1.15}

\setlength{\parindent}{1.15em}
\setlength{\parskip}{0.18em}
\raggedbottom

% ============================================================
% Colores: paleta original de inspiración académica británica
% ============================================================

\usepackage{xcolor}

\definecolor{AcademicNavy}{HTML}{17324D}
\definecolor{AcademicTeal}{HTML}{6F9390}
\definecolor{AcademicSlate}{HTML}{4A5561}
\definecolor{AcademicMist}{HTML}{EEF2F3}
\definecolor{AcademicRule}{HTML}{C9D1D6}

% ============================================================
% Matemáticas, tablas y figuras
% ============================================================

\usepackage{amsmath}
\usepackage{amssymb}
\usepackage{mathtools}

\usepackage{booktabs}
\usepackage{tabularx}
\usepackage{longtable}
\usepackage{array}
\usepackage{multirow}
\usepackage{threeparttable}

\usepackage{graphicx}
\usepackage{float}
\usepackage{subcaption}
\usepackage[section]{placeins}

\graphicspath{
    {../resultados/figuras/}
}

\usepackage{siunitx}
\sisetup{
    output-decimal-marker = {,},
    group-separator = {.},
    group-minimum-digits = 4
}

% ============================================================
% Títulos, listas, pies y navegación
% ============================================================

\usepackage{titlesec}
\usepackage{enumitem}
\usepackage{caption}
\usepackage{fancyhdr}

\titleformat{\section}
    {\Large\sffamily\bfseries\color{AcademicNavy}}
    {\thesection}
    {0.75em}
    {}

\titleformat{\subsection}
    {\large\sffamily\bfseries\color{AcademicSlate}}
    {\thesubsection}
    {0.70em}
    {}

\titleformat{\subsubsection}
    {\normalsize\sffamily\bfseries\color{AcademicSlate}}
    {\thesubsubsection}
    {0.65em}
    {}

\titlespacing*{\section}{0pt}{2.2ex plus 0.7ex minus 0.3ex}{1.2ex}
\titlespacing*{\subsection}{0pt}{1.8ex plus 0.5ex minus 0.2ex}{0.8ex}

\setlist[itemize]{
    leftmargin=1.6em,
    itemsep=0.25em,
    topsep=0.35em
}

\setlist[enumerate]{
    leftmargin=1.8em,
    itemsep=0.25em,
    topsep=0.35em
}

\captionsetup{
    font=small,
    labelfont={bf,sf,color=AcademicNavy},
    textfont=small,
    justification=justified,
    singlelinecheck=false,
    skip=7pt
}

\pagestyle{fancy}
\fancyhf{}
\fancyhead[L]{\sffamily\small MCC002 \textperiodcentered\ Grupo 5}
\fancyhead[R]{\sffamily\small\nouppercase{\leftmark}}
\fancyfoot[C]{\sffamily\small\thepage}

\renewcommand{\headrulewidth}{0.35pt}
\renewcommand{\footrulewidth}{0pt}

% ============================================================
% Bibliografía e hipervínculos
% ============================================================

\usepackage{csquotes}
\usepackage[
    backend=biber,
    style=apa,
    sorting=nyt,
    maxcitenames=2,
    maxbibnames=99
]{biblatex}

\usepackage[
    colorlinks=true,
    linkcolor=AcademicNavy,
    citecolor=AcademicNavy,
    urlcolor=AcademicTeal,
    pdfborder={0 0 0}
]{hyperref}

\hypersetup{
    pdftitle={Informe final MCC002 - Grupo 5},
    pdfauthor={Armando Castro Chaupis, Henry Sánchez Alvarado y Alex Segura Núñez},
    pdfsubject={Riesgo crediticio y predicción cuantitativa},
    pdfkeywords={riesgo crediticio, regresión logística, inferencia, Bayes, German Credit Data}
}

% ============================================================
% Entornos propios
% ============================================================

\newenvironment{hallazgo}
{
    \begin{quote}
    \color{AcademicNavy}
    \sffamily
    \textbf{Hallazgo principal.}\;
}
{
    \end{quote}
}

\newcommand{\variable}[1]{\texttt{#1}}
\newcommand{\pvalor}{\ensuremath{p\text{-valor}}}
\newcommand{\AUC}{\ensuremath{\mathrm{AUC}}}
\newcommand{\OR}{\ensuremath{\mathrm{OR}}}

% Control tipográfico
\widowpenalty=10000
\clubpenalty=10000
\displaywidowpenalty=10000
\emergencystretch=2em
